\section*{अध्याय \ref{ch:g3:problems} --- समस्याएँ हल करना}

\begin{solution}{exo:g3:problems:1}
झुंडों का मिलना: जोड़। चेरी बाँटना: भाग। पेंसिलों के डिब्बे:
गुणा। कार्ड दे देना: घटाव।
\end{solution}

\begin{solution}{exo:g3:problems:2}
पूछा गया: अब कितनी कारें। पता है: $246$, फिर $158$ और।
संक्रिया: जोड़। गणना: $246 + 158 = 404$। वाक्य: पार्किंग में अब
$404$ कारें हैं।
\end{solution}

\begin{solution}{exo:g3:problems:3}
समूह बनाने वाला भाग: $90 \div 9 = 10$। रस्सी से $10$ टुकड़े बनते हैं।
\end{solution}

\begin{solution}{exo:g3:problems:4}
गुणा: $26 \times 8 = 208$। टिकटों के $208$ लगेंगे।
\end{solution}

\begin{solution}{exo:g3:problems:5}
तुलना वाली पट्टी: एम्मा की पट्टी $64$, ह्यूगो की $29$ छोटी।
घटाव: $64 - 29 = 35$। ह्यूगो के पास $35$ स्टिकर हैं।
\end{solution}

\begin{solution}{exo:g3:problems:6}
मिले: $6 \times 45 = 270$ सेब। बचे:
$270 - 178 = 92$ सेब।
\end{solution}

\begin{solution}{exo:g3:problems:7}
विज्ञान की किताबें: $6 + 8 = 14$। उपन्यास: सप्ताह 2 में $12 - 9 = 3$ ज़्यादा।
\end{solution}

\begin{solution}{exo:g3:problems:8}
सोमवार $12$, मंगलवार $20$, बुधवार $12$: कुल
$12 + 20 + 12 = 44$ केक। सबसे अच्छा दिन मंगलवार था।
\end{solution}

\begin{solution}{exo:g3:problems:9}
बेकार संख्या क़ीमत ($12$) है। भाग:
$210 \div 35 = 6$ (क्योंकि $35 \times 6 = 210$)। लियो को $6$ दिन चाहिए।
\end{solution}

\begin{solution}{exo:g3:problems:10}
बड़े: $2 \times 9 = 18$। बच्चे: $3 \times 4 = 12$। कुल क़ीमत:
$18 + 12 = 30$। बाक़ी: $50 - 30 = 20$। उन्हें $20$ वापस मिलते हैं।
\end{solution}

\begin{solution}{exo:g3:problems:11}
कई सही उत्तर हो सकते हैं --- जैसे: ``एक फूलवाला $15$-$15$ फूलों
के $8$ गुच्छे बनाता है, फिर $20$ मुरझाए फूल निकाल देता है। कितने
फूल बचे?'' हल: $8 \times 15 = 120$, फिर $120 - 20 = 100$।

\end{solution}
